

Now, as we build up our multi-valley wavefunctions we consider what we will
describe here as the \emph{longitudinal single-axis basis wavefunction}. We
found that, for $\left(\kappa,i\right) = \left(A_{n},0\right)$,
$\left(E,i\right)$, and $\left(T_{n},z\right)$ the part of the wavefunction
lying along the $z$-axial Bloch function have the following forms
\begin{align}
    \psi_{\parallel_{0},m_{e}lk\sigma}^{\left(\kappa,z\right)}
        & \coloneqq
        \Phi_{\parallel_{0},m_{e}lk\sigma}^{\left(\kappa,z\right)}
        \eta^{\left(z\right)}_{\Omega^{\left(\kappa\right)}_{m_{e}l\sigma}}
\tag{\ref{single-axis:longitudinal:z:0}a}
,\\
    \psi_{\parallel_{1},mlk\sigma}^{\left(\kappa,z\right)}
        & \coloneqq
        \Phi_{\parallel_{1},mlk\sigma}^{\left(\kappa,z\right)}
        \eta^{\left(z\right)}_{\omega^{\left(\kappa\right)}\left(-1\right)^{m}}
\tag{\ref{single-axis:longitudinal:z:1}a}
.
\end{align}
The envelope basis functions associated with the basis wavefunctions above are constructed from our naive envelope basis functions and go as follows
\begin{align}
    \Phi_{\parallel_{0},m_{e}lk\sigma}^{\left(\kappa,z\right)}
        &=
            \phi_{klm_{e}\sigma;\left(-1\right)^{m_{e}/2}}^{(z)}
,\\
    \Phi_{\parallel_{1},mlk\sigma}^{\left(\kappa,z\right)}
        &=
          \phi_{klm\sigma;\left(-1\right)^{l}}^{(x)}
          +
          \omega^{\left(\kappa\right)}\Omega^{\left(\kappa\right)}_{ml\sigma}
          \phi_{klm\sigma;\left(-1\right)^{l+m}}^{(y)}
\tag{\ref{env:long:z:1}a}
.
\end{align}
Above we introduced the following constant.
\begin{equation}
\Omega^{\left(\kappa\right)}_{ml\sigma}
    \coloneqq
        \sigma\omega^{\left(\kappa\right)}
        \left(-1\right)^{l+m}\im^{+m}.
\end{equation}








A quick word on a notational brevity introduced above. We denote the two basis
states above by the subscripts $0$ and $1$. Each index on the wavefunction
belongs to that wavefunction, and should properly be subscripted by this label.
Since the author suspects that will introduce an untenable amount of cognitive
load within the reader, we use the convention that terms or tuples surrounded by
parentheses subscripted by $0$ and $1$ individually or $01$ together are
understood to have indices that are likewise subscripted by that index. As an
example, consider below.
\begin{equation}
\left(
    \tilde{\Psi}
    _{\varrho,mlk\xi}^{\left(\kappa,q\right)\left(\mathfrak{u}\right)}
\right)_{0}
    =
        \tilde{\Psi}
        _{
            \varrho_{0},m_{0}l_{0}k_{0}
            \xi_{0}}^{\left(\kappa_{0},q_{0}\right)\left(\mathfrak{u}_{0}\right)
        }
\end{equation}
We also extend this convention to superscripted or subscripted indices of
functions, namely
\begin{equation}
\tilde{\mathcal{V}}_{\parallel\left\{\bm{r}\right\}}
^{\left(\kappa,i\right)_{01}}
    =
        \tilde{\mathcal{V}}_{\parallel\left\{\bm{r}\right\}}
        ^{\left(\kappa_{0},i_{0},\kappa_{1},i_{1}\right)}
\end{equation}
This also serves as an example to show how parenthesis subscripted by $01$
should be interpreted. Finally, terms surrounded by brackets subscripted by $01$
are understood to be repeated, first with the subscript $0$ and then with
subscript $1$. Often, this means the repeated symbols inside the bracket are to
be multiplied together, for example
\begin{equation}
\left[
    \tilde{\Psi}
    _{\varrho,mlk\xi}^{\left(\kappa,q\right)\left(\mathfrak{u}\right)}
\right]_{01}
    =
        \tilde{\Psi}
        _{
            \varrho_{0},m_{0}l_{0}k_{0}
            \xi_{0}}
                ^{\left(\kappa_{0},q_{0}\right)\left(\mathfrak{u}_{0}\right)
        }
        \tilde{\Psi}
        _{
            \varrho_{1},m_{1}l_{1}k_{1}
            \xi_{1}}
                ^{\left(\kappa_{1},q_{1}\right)\left(\mathfrak{u}_{1}\right)
        }
.
\end{equation}

In the interests of making use of this software as flexible as possible, we
employ a potential of the form
\begin{equation}
\tilde{\mathcal{U}}_{\left\{\bm{r}\right\}}
    \coloneqq
        \mathfrak{E}^{\mathfrak{s}}
        \sum_{i}
        \left(\frac{r}{\mathfrak{r}}\right)^{\mathfrak{d}_{i}}
        \mathcal{U}_{i\left\{\Omega\right\}}
        \e^{
            -\tilde{\alpha}_{i\left\{\Omega\right\}}
            \frac{r}{\mathfrak{r}}
        }
        \e^{
            -\beta_{i\left\{\Omega\right\}}
            \left(\frac{r}{\mathfrak{r}}\right)^{2}
        }
\label{definition:potential:general}
\end{equation}
Above we introduce the integration "variable" $\left\{\Omega\right\}$. This
symbol denotes a solid angle functional dependence and we consider it
interchangeable with an expression using our angular spherical coordinates,
namely that $\left\{\Omega\right\} = \left\{\chi\phi\right\}$. As well,
$\mathfrak{E}$ represents the energy scale of our system and $\mathfrak{r}$
represents the length scale. Both have been defined in Appendix
\ref{appendix:units}. While we described the expression above as a "potential",
we see that by judicious choice of the constants (such as $\mathfrak{s}$,
$\mathfrak{d}$ and $\mathfrak{U}$) and functions (such as $\mathcal{U}$,
$\tilde{\alpha}$ and $\beta$) above this expression can be used for both the
overlap and Hamiltonian squared matrix elements in addition to the regular
Hamiltonian elements. We consider the potential above to be "complex" and apply
a tilde because while we expect the overall potential to be nondispersive and
real it is possible for individual terms to have complex coefficients and
functional dependence, depending upon how that potential is formulated.






We can start to simplify our integration by taking advantage of the cyclic
nature of the single-axis wavefunction. To wit, we may expand the full
wavefunctions in the matrix element formula and group by similar pairs of
valley-indices. Further we may exploit the transformation properties of the
single-axis wavefunction by using the cyclic transformations as a change of
integration variables to show that several terms of the integrand may be "folded
back over" each other. To that end, we may write
\begin{equation}
\left(\mathcal{M}\right)_{01}
    =
        \left(\mathcal{M}\right)_{\parallel,01}
            +
        \left(\mathcal{M}\right)_{\perp,01}
\end{equation}
where
\begin{equation}
\left(\mathcal{M}\right)_{\parallel,01}
    =
        \int_{-\pi}^{+\pi}\diff\phi
        \int_{-1}^{+1}\diff\chi
        \int_{0}^{\infty} r^2\diff{r}
        \,
        \left(
            \tilde{\psi}
            _{\varrho,mlk\xi}^{\left(\kappa,z\right)\left(\mathfrak{u}\right)}
        \right)_{0}^{*}
        \tilde{\mathcal{V}}_{\parallel\left\{\bm{r}\right\}}
        ^{\left(\kappa,i\right)_{01}}
        \left(
            \tilde{\psi}
            _{\varrho,mlk\xi}^{\left(\kappa,z\right)\left(\mathfrak{u}\right)}
        \right)_{1},
\end{equation}
and
\begin{equation}
\left(\mathcal{M}\right)_{\perp,01}
    =
        \int_{-\pi}^{+\pi}\diff\phi
        \int_{-1}^{+1}\diff\chi
        \int_{0}^{\infty} r^2\diff{r}
        \,
        \left(
            \tilde{\psi}
            _{\varrho,mlk\xi}^{\left(\kappa,y\right)\left(\mathfrak{u}\right)}
        \right)_{0}^{*}
        \tilde{\mathcal{V}}_{\perp\left\{\bm{r}\right\}}
        ^{\left(\kappa,i\right)_{01}}
        \left(
            \tilde{\psi}
            _{\varrho,mlk\xi}^{\left(\kappa,x\right)\left(\mathfrak{u}\right)}
        \right)_{1} .       
\end{equation}
The expressions denoted by $\tilde{\mathcal{V}}$ are the original general
potential form "folded over" together after change of integration variables
using the $T_{d}$ symmetry operations. These new forms may be compactly
expressed as below.
\begin{equation}
\tilde{\mathcal{V}}_{\parallel\left\{\bm{r}\right\}}
^{\left(\kappa,i\right)_{01}}
    \coloneqq
        \left\{\bm{\alpha}^{\left(\kappa,i\right)_{0}}\right\}^{\dagger}
        \begin{pmatrix}
            \tilde{\mathcal{U}}_{\left\{xyz\right\}} & 0 & 0 \\
            0 & \tilde{\mathcal{U}}_{\left\{zxy\right\}} & 0 \\
            0 & 0 & \tilde{\mathcal{U}}_{\left\{yzx\right\}} \\
        \end{pmatrix}
        \bm{\alpha}^{\left(\kappa,i\right)_{1}}
,
\end{equation}
and
\begin{equation}
\tilde{\mathcal{V}}_{\perp\left\{\bm{r}\right\}}
^{\left(\kappa,i\right)_{01}}
    \coloneqq
        \left\{\bm{\alpha}^{\left(\kappa,i\right)_{0}}\right\}^{\dagger}
        \begin{pmatrix}
            0
            & \xi_{0}\xi_{1}\tilde{\mathcal{U}}_{\left\{xzy\right\}}
            & \tilde{\mathcal{U}}_{\left\{zxy\right\}}
        \\
            \tilde{\mathcal{U}}_{\left\{yzx\right\}}
            & 0
            & \xi_{0}\xi_{1}\tilde{\mathcal{U}}_{\left\{yxz\right\}}
        \\
            \xi_{0}\xi_{1}\tilde{\mathcal{U}}_{\left\{zyx\right\}}
            & \tilde{\mathcal{U}}_{\left\{xyz\right\}}
            & 0
        \\
        \end{pmatrix}
        \bm{\alpha}^{\left(\kappa,i\right)_{1}}
,
\end{equation}
We note that we haven't a priori set $\kappa_{0} = \kappa_{1}$, because we wish
to examine and include, in addition to the potential of the system, other
external potentials such as applied electric fields that might break the
symmetry and allowing mixing of states with different $\kappa$. For our purposes
here, we note that while a general potential of the form we gave for
$\tilde{\mathcal{U}}$ does not have any particular sense of symmetry, it is well
known that general functions can be written as a sum of functions of definite
symmetry. To wit,
\begin{equation}
\tilde{\mathcal{U}}
    =
    \tilde{\mathcal{U}}^{\left(A_{1}\right)}
        +
    \tilde{\mathcal{U}}^{\left(A_{2}\right)}
        +
    \tilde{\mathcal{U}}^{\left(E,0\right)}
        +
    \tilde{\mathcal{U}}^{\left(E,1\right)}
        +
    \sum_{q=x,y,z}
    \left(
        \tilde{\mathcal{U}}^{\left(T_{1},q\right)}
            +
        \tilde{\mathcal{U}}^{\left(T_{2},q\right)}
    \right)
\end{equation}
We examine such selection rules and how they might simplify the above
expressions for $\tilde{\mathcal{V}}$ in the
\href{https://gitlab.com/cantankerous-impurity/Documentation/-/tree/core/SUNY--Buffalo/Dissertation/Supplementary%20Materials?ref_type=heads}{Supplementary Materials}. In any case, for this
derivation we only consider potentials of a given symmetry. The process of
resolving a general potential into symmetric components doesn't significantly
alter the underlying summation structure of the potential we defined in
\eqref{definition:potential:general} and so we will continue to use it. To that
end, we write
\begin{equation}
\tilde{\mathcal{V}}_{\aleph\left\{\bm{r}\right\}}
^{\left(\kappa^{\prime},j\right)\left(\kappa,i\right)_{01}}
    =
        \mathfrak{E}^{\mathfrak{s}}
        \sum_{i}
            \left(\frac{r}{\mathfrak{r}}\right)^{\mathfrak{d}_{i}}
            \mathcal{V}_{i\left\{\chi,\phi\right\}}
            \e^{
                -\tilde{\alpha}_{i\left\{\chi,\phi\right\}}
                \frac{r}{\mathfrak{r}}
            }
            \e^{
                -\beta_{i\left\{\chi,\phi\right\}}
                \left(\frac{r}{\mathfrak{r}}\right)^{2}
            }
\label{definition:potential:symmetrize}
.
\end{equation}

We may employ this form within our integration formula. One of the computational
benefits of the potential form we have chosen is that we may complete the radial
integration analytically, which we have seen will simplify the numerical
computational complexity of the integration over the solid angle degrees of
freedom.
\begin{align}
\left(\mathcal{M}\right)_{\aleph,01}^{\left(\varkappa,j\right)}
    &=
        \int_{-\pi}^{+\pi}\diff\phi
        \int_{-1}^{+1}\diff\chi
        \int_{0}^{\infty} r^2\diff{r}
        \,
        \xi_{0}
        \left[
            \tilde{\psi}
            _{
                \varrho,mlk\xi}^{\left(\kappa,q_{\aleph}\right)
                \left(\mathfrak{u}\right)
            }
        \right]_{01}
        \tilde{\mathcal{V}}_{\aleph\left\{\bm{r}\right\}}
        ^{\left(\kappa^{\prime},j\right)\left(\kappa,i\right)_{01}}
,\\
    &=
        \int_{\mathcal{S}_{0}}\diff\Omega
        \,
        \sum_{i}
        \sum_{s_{0}=-\mathfrak{u}_{0}}^{s_{0}=+\mathfrak{u}_{0}}
        \sum_{s_{1}=-\mathfrak{u}_{1}}^{s_{1}=+\mathfrak{u}_{1}}
        \mathfrak{C}
            _{i\left(ml;\eta\lambda,s\right)_{01}}
            ^{\left(\mathfrak{u}\right)_{01}}
        \mathcal{I}_{i\left(s\right)_{01}\left\{\Omega\right\}}
,\\
\mathfrak{C}
    _{i\left(ml;\eta\lambda,s\right)_{01}}
    ^{\left(\mathfrak{u}\right)_{01}}
    &\coloneqq
        \xi_{0}
        \mathfrak{E}^{\mathfrak{u}_{0}+\mathfrak{u}_{1}+\mathfrak{s}}
        \mathfrak{r}^{-\mathfrak{d}_{i}}
        \left[
            \frac
            { k!\mathfrak{H}_{ml;\lambda,s}^{\left(\mathfrak{u}\right)} }
            % { \mathfrak{K}!\left(\eta/\mathfrak{r}\right)^{2\mathfrak{u}} }
            {
                \mathfrak{K}!
                \left(\frac{\eta}{\mathfrak{r}}\right)^{2\mathfrak{u}}
            }
        \right]_{01}
.
\end{align}
We will see $\mathfrak{K}$ emerge later in this appendix. For now, we merely
remark that
\begin{equation}
    \mathfrak{K}_{n} = \left(k + \mathfrak{u} + \delta_{s\bar{1}}\right)_{n}
.
\end{equation}
Because we expect that, in general, the integration over solid angle will be
conducted numerically, we do not exploit the linear nature of integration to
pull out the summation symbols. We find that for our integrals it is
computationally much simpler and faster to compute one numerical integration
for a summed-over integrand rather then complete many numerical integrations
that are then summed. We have introduced a few sequences which we will define
below,
\begin{align}
\bm{\mathcal{I}}_{i\left(s\right)_{01}\left\{\chi,\phi\right\}}
    &\coloneqq
        \left[
            P_{\left(l+2s\right)m;\lambda\left\{\chi_{q}\right\}}
            Q_{m\zeta\left\{\phi_{q}\right\}}
        \right]_{01}
        \mathcal{V}_{i\left\{\Omega\right\}}
        \mathcal{R}_{i\left(s\right)_{01}\left\{\Omega\right\}}
,\\
\bm{\mathcal{R}}_{i\left(s\right)_{01}\left\{\Omega\right\}}
    &\coloneqq
        \int_{0}^{\infty} \diff{r}\,
        r^{\mathfrak{d}_{i}+2}
        \left[
            \frac{\mathfrak{K}!}{k!}
            \left(F_{s}^{\prime}\right)
            _{lk;\eta\lambda\left\{r\chi_{q}\right\}}
            ^{\left(W\mathfrak{u}\right)}
            \Theta^{\left(q\right)}
        \right]_{01}
        \e^
        {
                -
            \tilde{\alpha}_{i\left\{\Omega\right\}}\frac{r}{\mathfrak{r}}
                -
            \beta_{i\left\{\Omega\right\}}
            \left(\frac{r}{\mathfrak{r}}\right)^{2}
        }
.
\end{align}
We might note here that we have not explicitly included the Bloch functions in
our formulae above. However, with this form, incorporation of the Bloch
functions is fair straightforward. More details on how the full Bloch functions
may be incorporated into these formulae without much additional complexity may
be found in the
\href{https://gitlab.com/cantankerous-impurity/Documentation/-/tree/core/SUNY--Buffalo/Dissertation/Supplementary%20Materials?ref_type=heads}{Supplementary Materials}.

We recall that 
\begin{align}
\left[
    \left(F_{s}^{\prime}\right)
    _{lk;\eta\lambda\left\{r\chi_{q}\right\}}
    ^{\left(W\mathfrak{u}\right)}
\right]_{01}
    &=
        \left[
            \left(\sqrt{\frac{\vartheta}{\eta}}\right)^{3}
            \sqrt{W}
            \rho^{
                \frac{l}{W}
                -\mathfrak{u}\delta_{W1}-\delta_{s\bar{1}}
            }
            \e^{-\frac{\rho}{2}}
        \right]_{01}
        \mathcal{L}_{\left(s\left\{\rho\right\}\right)_{01}}
,\\
    &=
        r^{-3}W
        \left[
            \rho^{
                \frac{l+\frac{3}{2}}{W}
                -\mathfrak{u}\delta_{W1}-\delta_{s\bar{1}}
            }
            \e^{-\frac{\rho}{2}}
        \right]_{01}
        \mathcal{L}_{\left(s\left\{\rho\right\}\right)_{01}}
,\\
\mathcal{L}_{\left(s\left\{\rho\right\}\right)_{01}}
    &\coloneqq
        \left[
            \frac{\mathfrak{K}!}{k!}
            \sum_{t=-2\mathfrak{u}}^{\mathfrak{u}+\delta_{s\bar{1}}}
            h^{\left(\mathfrak{u}\right)}_{kl,st}
            L
            _{\left(k+t\right)\left\{\rho\right\}}
            ^{\left(\mathfrak{l}+\mathfrak{u}-\delta_{s\bar{1}}\right)}
        \right]_{01}
.   
\end{align}
We may make our notation even further compact. If we introduce
\begin{equation}
G_{\left\{r\right\}\left(s\left\{\chi_{q}\right\}\right)_{01}}
    \coloneqq
        W
        r^{\mathfrak{d}_{i}-1}
        \left[
            \rho^{
                \frac{l+\frac{3}{2}}{W}
                -\mathfrak{u}\delta_{W1}-\delta_{s\bar{1}}
            }
            \e^{-\frac{\rho}{2}}
        \right]_{01}
,
\end{equation}
than we may write our analytic integral as
\begin{equation}
\mathcal{R}_{i\left\{\Omega\right\}}
    =
        \int_{0}^{\infty} \diff{r}\,
        G_{\left\{r\right\}\left(s\left\{\chi_{q}\right\}\right)_{01}}
        \mathcal{L}_{\left(s\left\{\rho\right\}\right)_{01}}
        \e^
        {
                -
            \tilde{\alpha}_{i\left\{\Omega\right\}}
            \frac{r}{\mathfrak{r}}
                -
            \beta_{i\left\{\Omega\right\}}
            \left(\frac{r}{\mathfrak{r}}\right)^{2}
        }
.
\end{equation}
We do not assume that the anisotropy of the two states are the same in general,
and so $\rho_{0} \not\equiv \rho_{1}$ over the integration domain. We introduce
a new "averaged" variable. We define
\begin{equation}
\hat{\rho}_{\left\{\Omega\right\}}
    \coloneqq \frac{\rho_{0} + \rho_{1}}{2}
\end{equation}
which yields
\begin{equation}
\left[
    \e^{-\frac{\rho}{2}}
\right]_{01}
    =
        \e^{-\hat{\rho}}.
\end{equation}
We find it convenient to introduce a new radial scaling $\check{r}$ with solid
angle dependence so that we may write our "averaged" $\hat{\rho}$ as
\begin{equation}
\hat{\rho}_{\left\{\Omega\right\}}
    =
        \left(\frac{r}{\check{r}_{\left\{\Omega\right\}}}\right)^{W}
.
\label{definition:hat-rho}
\end{equation}
We will define $\check{r}$ below. If we employ this form of $\hat{\rho}$ we can
define our original $\rho_{n}$ in the following terms below
\begin{align}
\Gamma_{n\left\{\Omega\right\}}
    &\coloneqq
        \sqrt{
            \check{r}_{\left\{\Omega\right\}}
            \left(\frac{\vartheta_{\left\{\Omega\right\}}}{\eta}\right)_{n}
        },
\\
\rho_{n}
    &=
        \Gamma_{n\left\{\Omega\right\}}^{2W}\hat{\rho}
.
\end{align}
Using these new expressions we can simplify and rewrite part of our integrand
using $\hat{\rho}$ as
\begin{align}
\left[
    \rho^{
        \frac{l+\frac{3}{2}}{W}
        -\mathfrak{u}\delta_{W1}-\delta_{s\bar{1}}
    }
\right]_{01}
   & =
        \left[
            \Gamma
            _{\left\{\Omega\right\}}
            ^{\mathfrak{L}-2W\delta_{s\bar{1}}}
        \right]_{01}
        \hat{\rho}^{
            \frac{\mathfrak{L}_{0}+\mathfrak{L}_{1}}{2W}
            -\delta_{s_{0}\bar{1}}-\delta_{s_{1}\bar{1}}
        }
,\\
\mathfrak{L}_{n}
    &\coloneqq
        2l_{n} + 3 - 2\mathfrak{u}_{n}\delta_{W1}
.
\end{align}

Through experience, we have found it convenient in the development of our
software to rewrite $\rho$ in terms of some accessory constants. We have an effective radial functions that goes as
\begin{equation}
\check{r}_{\left\{\Omega\right\}}
    \coloneqq
        \frac
        {\hat{\eta}}
        {\hat{\vartheta}_{\left\{\Omega\right\}}\Xi_{\left\{\Omega\right\}}}
\end{equation}
where we defined an additional parameter
\begin{equation}
\hat{\eta}^{W}
    \coloneqq
        \frac
        { 2\left(\eta_{0}\eta_{1}\right)^{W} }
        { \eta_{0}^{W} + \eta_{1}^{W} }
\end{equation}
and some additional functions
\begin{align}
\hat{\vartheta}_{\left\{\Omega\right\}}^{W}
    &\coloneqq
        \frac{
            \vartheta_{0\left\{\Omega\right\}}^{W}
                +
            \vartheta_{1\left\{\Omega\right\}}^{W}
        }
        {2}
\\
\Xi_{\left\{\Omega\right\}}^{W}
    &\coloneqq
        1 + \mathfrak{f}\epsilon_{\left\{\Omega\right\}}
\end{align}
The $\Xi$ function depends on the mismatch between the parameters of the states,
which go as
\begin{align}
\mathfrak{f}
    & \coloneqq 
        \frac{\eta_{1}^{W}-\eta_{0}^{W}}{\eta_{0}^{W}+\eta_{1}^{W}}
\\
\epsilon_{\left\{\Omega\right\}}
    & \coloneqq 
        \frac
        {\vartheta_{0}^{W}-\vartheta_{1}^{W}}
        {\vartheta_{0}^{W}+\vartheta_{1}^{W}}
\end{align}
Within a subbasis $\eta_{0} = \eta_{1}$ and $\lambda_{0} = \lambda_{1}$ so many
of the above expressions simplify. However, $\epsilon \ne 0$ can emerge when
integrating over matrix elements involving single-valley states of different
axes.

With these accessor functions and parameters, we can rewrite one of the factors
of our integrand as such
\begin{align}
G_{\left\{r\right\}\left(s\left\{\chi_{q}\right\}\right)_{01}}
    &=
        W
        r^{\mathfrak{d}_{i}-1}
        \left[
            \rho^{
                \frac{l+\frac{3}{2}}{W}
                -\mathfrak{u}\delta_{W1}-\delta_{s\bar{1}}
            }
            \e^{-\frac{\rho}{2}}
        \right]_{01}
,\\
    &=
        W
        \left(
            \hat{\rho}^{\frac{1}{W}}\check{r}_{\left\{\Omega\right\}}
        \right)
        ^{\mathfrak{d}_{i}-1}
        \left[
            \Gamma
            _{\left\{\Omega\right\}}
            ^{\mathfrak{L}-2W\delta_{s\bar{1}}}
        \right]_{01}
        \hat{\rho}^{
            \frac{\mathfrak{L}_{0}+\mathfrak{L}_{1}}{2W}
            -\delta_{s_{0}\bar{1}}-\delta_{s_{1}\bar{1}}
        }
        \e^{-\hat{\rho}}
,\\
    &=
        \frac{W}{\check{r}_{\left\{\Omega\right\}}}
        \hat{\rho}^{1-\frac{1}{W}}
        H_{\left\{\hat{\rho}\right\}\left(s\left\{\chi_{q}\right\}\right)_{01}}
        \e^{-\hat{\rho}}
.
\end{align}
Above we have introduced a new function and some associated constants, which go
as
\begin{align}
H_{\left\{\hat{\rho}\right\}\left(s\left\{\chi_{q}\right\}\right)_{01}}
    &\coloneqq
        \check{r}_{\left\{\Omega\right\}}^{\mathfrak{d}_{i}}
        \left[
            \Gamma_{\left\{\Omega\right\}}^{\mathfrak{L}-2W\delta_{s\bar{1}}}
        \right]_{01}
        \hat{\rho}^{
            \frac{\hat{\mathfrak{L}}_{\mathfrak{d}_{i}}}{W}
            -\delta_{s_{0}\bar{1}}-\delta_{s_{1}\bar{1}}
            -1
        }
,\\
\hat{\mathfrak{L}}_{\mathfrak{d}}
    &\coloneqq
        \frac{\mathfrak{L}_{0}+\mathfrak{L}_{1}}{2} + \mathfrak{d}
    = 
        l_{0} + l_{1}
        - \delta_{W1}\left(\mathfrak{u}_{0}+\mathfrak{u}_{1}\right)
        + \mathfrak{d} + 3
.
\end{align}
We note that 
\begin{equation}
\diff\hat{\rho}
    =
        \frac{W}{\hat{r}}\hat{\rho}^{1-\frac{1}{W}}\diff r
\end{equation} as such we may write
\begin{equation}
\diff r\,
G_{\left\{r\right\}\left(s\left\{\chi_{q}\right\}\right)_{01}}
    =
\diff\hat{\rho}\,
H_{\left\{\hat{\rho}\right\}\left(s\left\{\chi_{q}\right\}\right)_{01}}
\end{equation}
The expressions above suggests we should redefine
$\hat{\rho}_{\left\{\Omega\right\}}$ as our integration variable. Over the
radial integration we can consider $\left\{\Omega\right\}$ to be constant, and
so we will drop this subscript from expressions that form the radial integrand.
We will restore this subscripts as needed, and once the analytical radial
integration is completed we will ensure these subscripts are included in the
numerical solid-angle integrand. With this in mind we may write our analytic
integral over $\hat{\rho}$ as
\begin{equation}
\mathcal{R}_{i\left(s\right)_{01}\left\{\Omega\right\}}
    =
        \int_{0}^{\infty} \diff{\hat{\rho}}\,
        H_{\left\{\hat{\rho}\right\}\left(s\right)_{01}}
        \mathcal{L}_{\left(s\left\{\rho\right\}\right)_{01}}
        \e^{-\mathcal{F}_{i\left\{\hat{\rho}\right\}}}
.
\end{equation}
Expanding $H$ as previously defined yields
\begin{equation}
\mathcal{R}_{i\left(s\right)_{01}\left\{\Omega\right\}}
    =
        \check{r}_{\left\{\Omega\right\}}^{\mathfrak{d}_{i}}
        \left[
            \Gamma_{\left\{\Omega\right\}}^{\mathfrak{L}-2W\delta_{s\bar{1}}}
        \right]_{01}
        \bm{\mathcal{R}}_{i\left(s\right)_{01}\left\{\Omega\right\}}
,
\end{equation}
with the radial integral partitioned as $\bm{\mathcal{R}}$. We'll focus on the
radial integral, namely  
\begin{align}
\bm{\mathcal{R}}_{i\left(s\right)_{01}\left\{\Omega\right\}}
    &\coloneqq
        \int_{0}^{\infty}\diff\hat{\rho}\,
        \hat{\rho}^{
            \frac{\hat{\mathfrak{L}}_{\mathfrak{d}_{i}}}{W}
            -\delta_{s_{0}\bar{1}}-\delta_{s_{1}\bar{1}}
            -1
        }
        \mathcal{L}_{\left(s\left\{\rho\right\}\right)_{01}}
        \e^{-\mathcal{F}_{i\left\{\hat{\rho}\right\}}}
,\\
\mathcal{L}_{\left(s\left\{\rho\right\}\right)_{01}}
    &\coloneqq
        \left[
            \frac{\mathfrak{K}!}{k!}
            \sum_{t=-2\mathfrak{u}}^{\mathfrak{u}+\delta_{s\bar{1}}}
            h^{\left(\mathfrak{u}\right)}_{kl,st}
            L
            _{\left(k+t\right)\left\{\rho\right\}}
            ^{\left(\mathfrak{l}+\mathfrak{u}-\delta_{s\bar{1}}\right)}
        \right]_{01}
\label{ex:sometimes_repeated_symbols_means_multiply}
.
\end{align}
In the analytical radial integral above we introduced a new function, 
\begin{align}
\mathcal{F}_{i\left\{\hat{\rho}\right\}}
    &\coloneqq
        \hat{\rho}
            +
        \tilde{\alpha}_{i}\left(\frac{\check{r}}{\mathfrak{r}}\right)
        \hat{\rho}^{\frac{1}{W}}
            +
        \beta_{i}\left(\frac{\check{r}}{\mathfrak{r}}\right)^{2}
        \hat{\rho}^{\frac{2}{W}}
,\\
    &=
        \left(1+\mathcal{A}_{i}\right)\hat{\rho}
            +
        \mathcal{B}_{i}\hat{\rho}^{\frac{2}{W^{2}}}
,
\end{align}
where 
\begin{align}
\mathcal{A}_{i}
    &\coloneqq
    \delta_{W1}
    \tilde{\alpha}_{i}
    \left(\frac{\check{r}}{\mathfrak{r}}\right)
        +
    \delta_{W2}
    \beta_{i}
    \left(\frac{\check{r}}{\mathfrak{r}}\right)^{2}
\label{def:A}
,\\
\mathcal{B}_{i}
    &\coloneqq
    \delta_{W1}
    \beta_{i}
    \left(\frac{\check{r}}{\mathfrak{r}}\right)^{2}
        +
    \delta_{W2}
    \tilde{\alpha}_{i}
    \left(\frac{\check{r}}{\mathfrak{r}}\right)
\label{def:B}
\end{align}
We can rewrite $\bm{\mathcal{R}}$ by pulling the summation over
$\left(t\right)_{01}$ outside the analytical radial integration.
\begin{align}
\bm{\mathcal{R}}_{i\left(s\right)_{01}}
    &=
        \left[
            \sum_{t=-2\mathfrak{u}}^{\mathfrak{u}+\delta_{s\bar{1}}}
            \frac{\mathfrak{K}!}{k!}
            h^{\left(\mathfrak{u}\right)}_{\left(kl,st\right)}
        \right]_{01}
        \mathcal{R}_{i\left(st\right)_{01}}^{\prime}
,\\
\mathcal{R}_{i\left(st\right)_{01}}^{\prime}
    &\coloneqq
        \int_{0}^{\infty}\diff\hat{\rho}\,
            \hat{\rho}^{
                \frac{\hat{\mathfrak{L}}_{\mathfrak{d}_{i}}}{W}
                -\delta_{s_{0}\bar{1}}-\delta_{s_{1}\bar{1}}
                -1
            }
            \left[
                L
                _{\left(k+v\right)\left\{\rho\right\}}
                ^{\left(\mathfrak{l}+\mathfrak{u}-\delta_{s\bar{1}}\right)}
            \right]_{01}
            \e^{-\mathcal{F}_{i\left\{\hat{\rho}\right\}}}
.
\end{align}
In analogy to our rules for subscripted brackets, we introduce a related rule
for summation symbols when they are the only symbols included in the brackets.
We want to be clear that this rule might be counterintuitive when first seen
because it represents an abuse of commonly understood rules of parenthesis and
bracket in arithmetic. However, we think the formal meaning of repeated symbols
is clear and it can be understood by the reader that in certain contexts
repeated symbols do necessarily mean "product" exclusively. As such, we want to
make it clear that symbols outside and to the right of subscripted brackets are
understood to be part of the summation, not outside it as might be expected.
Namely, we define the following.
\begin{equation}
\left[\sum_{n=a}^{b}\right]_{01}
\left(f_{n}\right)_{01}
    \coloneqq
        \sum_{n_{0}=a_{0}}^{b_{0}}
        \sum_{n_{1}=a_{1}}^{b_{1}}
        f_{n_{0}n_{1}}
\end{equation}
For example, equation \eqref{ex:first_use_of_bracket_summation} above can be
considered and understood in light of the definition above to be such that
\begin{equation}
\left[\sum_{s=-\mathfrak{u}}^{s=+\mathfrak{u}}\right]_{01}
\left(
    \mathfrak{C}
    _{i\left(ml;\eta\lambda,s\right)}
    ^{\left(\mathfrak{u}\right)}
    \mathcal{I}_{i\left(s\right)\left\{\Omega\right\}}^{\prime}
\right)_{01}
    =
        \sum_{s_{0}=-\mathfrak{u}_{0}}^{s_{0}=+\mathfrak{u}_{0}}
        \sum_{s_{1}=-\mathfrak{u}_{1}}^{s_{1}=+\mathfrak{u}_{1}}
        \mathfrak{C}
        _{i\left(ml;\eta\lambda,s\right)_{01}}
        ^{\left(\mathfrak{u}\right)_{01}}
        \mathcal{I}_{i\left(s\right)_{01}\left\{\Omega\right\}}^{\prime}
.
\end{equation}
While it represents a slight abuse of standard bracket notation, we hope that it
is clear that this notation has the benefit of compacting otherwise long and
repetitive expressions involving double or greater summation. To prevent
confusion, we note that forms with brackets containing summation symbol AND
summations terms, for example as in equation
\eqref{ex:sometimes_repeated_symbols_means_multiply}
should continue to be understood as a product of terms subscripted with $0$ and
$1$ respectively.

To complete the radial integration, we consider an integral of the form given
below.
\begin{align}
\mathcal{R}^{\prime}
&\coloneqq
    \int_{0}^{\infty}\diff\rho\,
        \rho^{\mathfrak{g}-1}
        L
        _{\mathfrak{T}_{0}\left\{\rho\left(1+\Delta\right)\right\}}
        ^{\left(\mathfrak{a}_{0}\right)}
        L
        _{\mathfrak{T}_{1}\left\{\rho\left(1-\Delta\right)\right\}}
        ^{\left(\mathfrak{a}_{1}\right)}
        \e^{
            -\left(1+\mathcal{A}\right)\rho
            -\mathcal{B}\rho^{\frac{2}{W^{2}}}
        }
,\\
&=
    \left[
        \sum_{u=0}^{\mathfrak{T}}
        \sum_{v=u}^{\mathfrak{T}}
    \right]_{01}
    \mathfrak{F}_{\mathfrak{g}\left(uv\right)_{01}}
    \mathcal{G}
    _{
        \left(W\mathfrak{g}\right);
        \left(u\right)_{01}\left\{\Delta\mathcal{Z}_{\mathcal{G}}\right\}
    }
    ^{\left(W\right)}
    \mathcal{H}
    _{\left(W\mathfrak{g}+Ww\right)\left\{\mathcal{Z}_{\mathcal{H}}\right\}}
    ^{\left(W\right)}
.
\end{align}

The $L$ functions used above are the associated Laguerre functions. The
significance of the above expression is that it includes the coefficients of the
Laguerre polynomials. It should be noted that the summations above can involve
sequences of terms of near magnitude but alternating sign. This is a computation
ill-posed to be completed numerically, so we will endeavor to simplify the
expressions as much as possible. Some of the expressions used in the formula
above are defined below
\begin{align}
\mathfrak{F}_{\mathfrak{g}\left(uv\right)_{01}}
    &\coloneqq
        \left[
            \frac
            {
                \left(-1\right)^{u+v}
                \sqrt{
                    \mathfrak{T}!
                    \left(\mathfrak{T}+\mathfrak{a}\right)!
                }
            }
            {
                \left(\mathfrak{T}-v\right)!
                \left(\mathfrak{a}+v\right)!
                u!
                \left(v-u\right)!
            }
        \right]_{01}
        \Gamma\left(\mathfrak{g}+v_{0}+v_{1}\right)
,\\
\mathcal{G}
_{L;\left(u\right)_{01}\left\{\Delta\mathcal{Z}\right\}}
^{\left(W\right)}
    &\coloneqq
        \left(1-\mathcal{Z}\left(1+\Delta\right)\right)^{u_{0}}
        \left(1-\mathcal{Z}\left(1-\Delta\right)\right)^{u_{1}}
        \mathcal{Z}^{\frac{L}{W}}
,\\
\mathcal{H}_{M\left\{\mathcal{Z}\right\}}^{\left(W\right)}
    &\coloneqq
        \left(
            \delta_{W1}
            \frac{\sqrt{\pi}}{\Gamma\left[\frac{M+1}{2}\right]}
            \mathcal{Z}^{M}
            +
            \delta_{W2}
        \right)
        h_{M\left\{\mathcal{Z}\right\}}
,\\
\mathcal{Z}_{\mathcal{G}} &\coloneqq \frac{1}{1+\mathcal{A}}
\label{def:Z0}
,\\
\mathcal{Z}_{\mathcal{H}}
    &\coloneqq
        \delta_{W1}
        \left(
            \frac{1+\mathcal{A}}{2\sqrt{\mathcal{B}}}
        \right)
            +
        \delta_{W2}
        \left(
            \frac{\mathcal{B}}{2\sqrt{1+\mathcal{A}}}
        \right)
        \label{def:Z1}
.
\end{align}
Tthe expression $h_{WL}$ refers to the hypergeometric function defined in the
\href{https://gitlab.com/cantankerous-impurity/Documentation/-/tree/core/SUNY--Buffalo/Dissertation/Supplementary%20Materials?ref_type=heads}{Supplementary Materials}.

We can use the above expressions if we set the constants to appropriate values.
As we are defining new quantities and functions, we've taken the liberty of
explicitly including the dependence over solid-angle in the final expression of
indirect integration variables. We define or set the following
\begin{align}
\mathfrak{h}_{\mathfrak{d}\left(s\right)_{01}}
    &\coloneqq
W\mathfrak{g}
    =
    \hat{\mathfrak{L}}_{\mathfrak{d}}
    -W\left(\delta_{s_{0}\bar{1}}+\delta_{s_{1}\bar{1}}\right)
,\\
\mathfrak{a}_{n}
    &\coloneqq
        \left(\mathfrak{l}+\mathfrak{u}-\delta_{s\bar{1}}\right)_{n}
,\\
\mathfrak{T}_{n}
    &\coloneqq
        \left(k+t\right)_{n}
,\\
\Delta
    &\coloneqq
        \frac{\rho_{0}-\rho_{1}}{2\hat{\rho}}
.
\end{align}
We can further refine $\Delta$.
\begin{align}
\Delta
    &=
        \frac{\Gamma_{0}^{2W}-\Gamma_{1}^{2W}}{2}
        =
        \frac{\check{r}^{W}}{2}
        \left[
            \left(\frac{\vartheta}{\eta}\right)_{0}^{W}
                -
            \left(\frac{\vartheta}{\eta}\right)_{1}^{W}
        \right]
,\\
    &=
    \frac{1}{2}
    \left(\frac{\hat{\eta}}{\hat{\vartheta}\Xi}\right)^{W}
    \left(
        \frac
        {
            \left(\eta_{1}\vartheta_{0}\right)^{W}
                -
            \left(\eta_{0}\vartheta_{1}\right)^{W}
        }
        {\left[\eta\right]_{01}^{W}}
    \right)
,\\
&=
    \frac{1}{2}
    \left(
        \frac
        {
            4\left[\eta\right]_{01}^{W}
        }
        {
            \left(\eta_{0}^{W} + \eta_{1}^{W}\right)
            \left( \vartheta_{0}^{W} + \vartheta_{1}^{W} \right)
            \Xi^{W}
        }
    \right)
    \left(
        \frac
        {
            \left(\eta_{1}\vartheta_{0}\right)^{W}
                -
            \left(\eta_{0}\vartheta_{1}\right)^{W}
        }
        {
            \left[\eta\right]_{01}^{W}
        }
    \right)
,\\
&=
    \frac
    {
        \left( \eta_{1}^{W} - \eta_{0}^{W} \right)
        \left( \vartheta_{0}^{W} + \vartheta_{1}^{W} \right)
            +
        \left( \eta_{0}^{W} + \eta_{1}^{W} \right)
        \left( \vartheta_{0}^{W} - \vartheta_{1}^{W} \right)
    }
    {
        \left( \vartheta_{0}^{W} + \vartheta_{1}^{W} \right)
        \left(\eta_{0}^{W} + \eta_{1}^{W}\right)
        \Xi^{W}
    }
.
\end{align}
Distributing the denominator through the sum in the numerator, and identifying
existing quantities, we may state our final expression for $\Delta$ to be
\begin{equation}
\Delta_{\left\{\Omega\right\}}
    =
        \frac
        {\mathfrak{f}+\epsilon_{\left\{\Omega\right\}}}
        {1+\mathfrak{f}\epsilon_{\left\{\Omega\right\}}}
\end{equation}

We must also consider $\mathcal{Z}_{\mathcal{G}}$ and
$\mathcal{Z}_{\mathcal{H}}$. We had defined these quantities in terms of
$\mathcal{A}$ and $\mathcal{B}$. Here we recall that we had previously
identified those quantities as indirect integration variables, namely
\begin{align}
\tilde{\mathcal{A}}_{i}
    &\coloneqq
        \delta_{W1}
        \tilde{\alpha}_{i}
        \left(\frac{\check{r}}{\mathfrak{r}}\right)
            +
        \delta_{W2}
        \beta_{i}
        \left(\frac{\check{r}}{\mathfrak{r}}\right)^{2}
\tag{\ref{def:A}a}
,\\
\tilde{\mathcal{B}}_{i}
    &\coloneqq
        \delta_{W1}
        \beta_{i}
        \left(\frac{\check{r}}{\mathfrak{r}}\right)^{2}
            +
        \delta_{W2}
        \tilde{\alpha}_{i}
        \left(\frac{\check{r}}{\mathfrak{r}}\right) 
\tag{\ref{def:B}a}
\end{align}
Using these expressions for $\mathcal{A}$ and $\mathcal{B}$, we can write the
particular $\mathcal{Z}_{\mathcal{G}}$ and $\mathcal{Z}_{\mathcal{H}}$ we wish
to use as functions over solid-angle. To wit,
\begin{align}
\tilde{\mathcal{Z}}_{i\mathcal{G}\left\{\Omega\right\}}
&=
    \frac
    {1}
    {
        1
            +
        \delta_{W1}
        \tilde{\alpha}_{i\left\{\Omega\right\}}
        \left(\frac{ \check{r}_{\left\{\Omega\right\}} }{ \mathfrak{r} }\right)
            +
        \delta_{W2}
        \beta_{i\left\{\Omega\right\}}
        \left(
            \frac{ \check{r}_{\left\{\Omega\right\}} }{ \mathfrak{r} }
        \right)^{2}
    }
,\\
\tilde{\mathcal{Z}}_{i\mathcal{H}\left\{\Omega\right\}}
&=
    \frac
    {
        \delta_{W1}
            +
        \tilde{\alpha}_{i\left\{\Omega\right\}}
        \left(\frac{ \check{r}_{\left\{\Omega\right\}} }{ \mathfrak{r} }\right)
    }
    {
        \sqrt{
            \delta_{W2}
                +
            \beta_{i\left\{\Omega\right\}}
            \left(
                \frac{ \check{r}_{\left\{\Omega\right\}} }{ \mathfrak{r} }
            \right)^{2}
        }
    }    
.
\end{align}
While we wanted to ensure the solid-angle dependence of $\Delta$ and the
$\mathcal{Z}$ were clear in the expressions above, we will resume dropping this
dependence in expressions while examining the analytic radial integration. In
light of the arguments above, we can write
\begin{align}
\bm{\mathcal{R}}_{i\left(s\right)_{01}}
    &=
        \left[
            \sum_{t=-2\mathfrak{u}}^{\mathfrak{u}+\delta_{s\bar{1}}}
            \frac{\mathfrak{K}!}{k!}
            h^{\left(\mathfrak{u}\right)}_{\left(kl,st\right)}
        \right]_{01}
        \mathcal{R}_{i\left(st\right)_{01}}^{\prime}
,\\
\mathcal{R}_{i\left(st\right)_{01}}^{\prime}
    &=
        \left[
            \sum_{u=0}^{k+t}
            \sum_{v=u}^{k+t}
        \right]_{01}
        \mathfrak{F}
        _{
            \left(
                \frac{\mathfrak{h}_{\mathfrak{d}_{i}\left(s\right)_{01}}}{W}
            \right)
            \left(uv\right)_{01}
        }
        \bm{\mathcal{G}}_{i\left(su\right)_{01}}
        \bm{\mathcal{H}}_{i\left(s\right)_{01}\left(v_{0}+v_{1}\right)}
.
\end{align}
where we found it convenient to introduce the following sequences
\begin{align}
\bm{\mathcal{G}}_{i\left(su\right)_{01}}
    &\coloneqq
        \mathcal{G}
        _{
            \mathfrak{h}_{\mathfrak{d}_{i}\left(s\right)_{01}};
            \left(u\right)_{01}
            \left\{\Delta\tilde{\mathcal{Z}}_{i\mathcal{G}}\right\}
        }
        ^{\left(W\right)}
\label{definition:G-array}
,\\
\bm{\mathcal{H}}_{i\left(s\right)_{01}w}
    &\coloneqq
        \mathcal{H}_{
            \left(
                \mathfrak{h}_{\mathfrak{d}_{i}\left(s\right)_{01}}
                +Ww
            \right)
            \left\{\tilde{\mathcal{Z}}_{i\mathcal{H}}\right\}
        }^{\left(W\right)}
\label{definition:H-array}
.
\end{align}

Earlier, we had pulled out the summation over $t$ to complete the radial
integration analytically. We can reorganize the terms of the double summation in
$\mathcal{R}$ to allow the summation over $t$ to be completed. We introduce a
new limit for summations that goes as
\begin{equation}
    \mathfrak{K}_{n} = \left(k + \mathfrak{u} + \delta_{s\bar{1}}\right)_{n}.
\end{equation}
We recognize that $u$ can range from $0$ to $\mathfrak{K}$, $v$ can range from
$u$ to $\mathfrak{K}$, and $t$ must always be larger then $v-k$. We reorganize
the order of our summations, and we pull out to the front summations over
$u_{0}$, and $u_{1}$. We may write
\begin{equation}
\bm{\mathcal{R}}_{i\left(s\right)_{01}}
    =
        \left[
            \sum_{u=0}^{\mathfrak{K}}
            \sum_{v=u}^{\mathfrak{K}}
        \right]_{01}
        \mathfrak{A}_{\mathfrak{d}_{i}\left(suv\right)_{01}}
        \bm{\mathcal{G}}_{i\left(su\right)_{01}}
        \bm{\mathcal{H}}_{i\left(s\right)_{01}\left(v_{0}+v_{1}\right)}
\end{equation}
where we introduced
\begin{equation}
\mathfrak{A}_{\mathfrak{d}\left(suv\right)_{01}}
    \coloneqq
        \left[
            \sum
            _{t=\text{Max}\left[v-k,-2\mathfrak{u}\right]}
            ^{\mathfrak{u}+\delta_{s\bar{1}}}
            \frac{\mathfrak{K}!}{k!}
            h^{\left(\mathfrak{u}\right)}_{\left(kl,st\right)}
        \right]_{01}
        \mathfrak{F}_{
            \left(\frac{\mathfrak{h}_{\mathfrak{d}}}{W}\right)
            \left(uv\right)_{01}
        }
\end{equation}
However, we can further simplify this from a four-index summation to a
three-index summation by introducing $w = v_{0}+v_{1}$. We do this by
introducing a summation over a Kronecker delta, going as
\begin{align}
\bm{\mathcal{R}}_{i\left(s\right)_{01}}
    &=
        \left[
            \sum_{u=0}^{\mathfrak{K}}
            \sum_{v=u}^{\mathfrak{K}}
        \right]_{01}
        \mathfrak{A}_{\mathfrak{d}_{i}\left(suv\right)_{01}}
        \bm{\mathcal{G}}_{i\left(su\right)_{01}}
        \sum_{w=u_{0}+u_{1}}^{\mathfrak{K}_{0}+\mathfrak{K}_{1}}
        \delta_{w,v_{0}+v_{1}}
        \bm{\mathcal{H}}_{i\left(s\right)_{01}w}
,\\
    &=
        \left[ \sum_{u=0}^{\mathfrak{K}} \right]_{01}
        \sum_{w=u_{0}+u_{1}}^{\mathfrak{K}_{0}+\mathfrak{K}_{1}}
        \mathfrak{A}_{\mathfrak{d}_{i}\left(su\right)_{01}w}
        \bm{\mathcal{G}}_{i\left(su\right)_{01}}
        \bm{\mathcal{H}}_{i\left(s\right)_{01}w}
\label{equation:radial-integral}
.
\end{align}
where we introduce a new constant
\begin{equation}
\mathfrak{A}_{\mathfrak{d}\left(su\right)_{01}w}
    \coloneqq
        \left[ \sum_{v=u}^{\mathfrak{K}} \right]_{01}
        \delta_{w,v_{0}+v_{1}}\mathfrak{A}_{\mathfrak{d}\left(suv\right)_{01}}
.
\end{equation}
$\mathfrak{A}$ is constant over integration, and in particular its value does not
depend on any runtime parameters, as long as the basis and operator/potential
specification occurs at or prior to compile time. It is a potentially
complicated expression, and we take some effort to simplify it and attempt to
find a form of the expression that may be computed in reusable pieces. We can
consider 
\begin{align}
\mathfrak{A}_{\mathfrak{d}\left(su\right)_{01}w}
    &=
        \left(-1\right)^{u_{0}+u_{1}+w}
        \Gamma\left[
            \frac{\mathfrak{h}_{\mathfrak{d}\left(s\right)_{01}}}{W}+w
        \right]
        \mathfrak{A}_{\left(su\right)_{01}w}^{\prime}
\label{appendix:elements:precalc-coeffs-0}
,\\
\mathfrak{A}_{\left(su\right)_{01}w}^{\prime}
    &=
        \left[\sum_{v=u}^{\mathfrak{K}}\right]_{01}
        \delta_{w,v_{0}+v_{1}}
        \left[\mathfrak{A}_{suv}^{\prime\prime}\right]_{01}
.
\end{align}
The product of coefficients above has factors that depend only on information
from a single state. We can write this single-state coefficient as so,
\begin{align}
\mathfrak{A}_{suv}^{\prime\prime}
    &=
        \frac{1}{\sqrt{k!\left(k+\mathfrak{l}\right)!}}
        \frac{\mathfrak{K}!}{u!\left(v-u\right)!\left(\mathfrak{K}-v\right)!}
        \frac
        {\left(k+\mathfrak{l}\right)!}
        {\left(\mathfrak{a}-2\mathfrak{u}+v\right)!}
        \mathfrak{G}_{sv}
,\\
    &=
        \frac{1}{\sqrt{k!\left(k+\mathfrak{l}\right)!}}
        \begin{pmatrix}\mathfrak{K}\\u\end{pmatrix}
        \begin{pmatrix}\mathfrak{K}-u\\\mathfrak{K}-v\end{pmatrix}
        \left(k+\mathfrak{l}\right)^{\underline{\mathfrak{K}-v}}
        \mathfrak{G}_{sv}
,\\
\mathfrak{G}_{sv}
    &=
        \frac
        {\left(\mathfrak{a}-2\mathfrak{u} + v\right)!}
        {\left(\mathfrak{a}+ v\right)!}
        \mathfrak{H}_{sv}^{\left(\mathfrak{u}\right)}
,\\
    &=
        \left(1-\mathfrak{u}\right)\mathfrak{H}_{sv}^{\left(0\right)}
            +
        \mathfrak{u}
        \frac
        { \mathfrak{H}_{sv}^{\left(1\right)} }
        {
            \left(\mathfrak{l}-\delta_{s\bar{1}}+v+1\right)
            \left(\mathfrak{l}-\delta_{s\bar{1}}+v\right) }
.
\end{align}
We have
\begin{equation}
\mathfrak{H}_{sv}^{\left(\mathfrak{u}\right)}
    =
        \sum
        _{t=\text{Max}\left[v-k,-2\mathfrak{u}\right]}
        ^{\mathfrak{u}+\delta_{s\bar{1}}}
        h^{\left(\mathfrak{u}\right)}_{kl,st}
        \frac
        {
            \left(\mathfrak{K}-v\right)!
            \sqrt{\left(k+t\right)!\left(k+t+\mathfrak{a}\right)!}
        }
        {\left(k+t-v\right)!\sqrt{k!\left(k+\mathfrak{l}\right)!}}
.
\end{equation}
We also find it convenient to explicitly the separate the $\mathfrak{u} = 0$
(kinetic-off) and the $\mathfrak{u} = 1$ (kinetic-on) forms somewhat. I note
that when the kinetic operator is "off" that $s = 0$ and that $v-k\le 0$. This
means
\begin{equation}
\mathfrak{H}_{sv}^{\left(0\right)} = 1
.
\end{equation}
We also have
\begin{equation}
    \mathfrak{H}_{sv}^{\left(1\right)}
        =
            \sum_{t=\text{Max}\left[v-k,-2\right]}^{1+\delta_{s\bar{1}}}
            g_{kl,st}
            \frac
            {
                \left(\mathfrak{K}-v\right)!
                \sqrt{
                    \left(k+t\right)!
                    \left(k+t+\mathfrak{l}+1-\delta_{s\bar{1}}\right)!
                }
            }
            {\left(k+t-v\right)!\sqrt{k!\left(k+\mathfrak{l}\right)!}}
.
\end{equation}
We note that we can treat this summation as a polynomial in $\mathfrak{K}-v$.
Let us define
\begin{align}
\mathfrak{H}^{\prime}_{sv}
    &=
        \sum
        _{t=0}
        ^{\text{Min}\left[\mathfrak{N},3+\delta_{s\bar{1}}\right]}
        v^{t}
        \mathfrak{G}_{st}
,\\
\mathfrak{G}_{st}
    &\coloneqq
        g_{kl,s\left(1+\delta_{s\bar{1}}-t\right)}
        \frac
        {
            \sqrt{
                \left(k+1+\delta_{s\bar{1}}-t\right)!
                \left(k+\mathfrak{l}+2-t\right)!
            }
        }
        {
            \sqrt{k!\left(k+\mathfrak{l}\right)!}
        }
.
\end{align}
We may pull together these realizations and write
\begin{align}
\mathfrak{A}_{su\left(\mathfrak{K}-v\right)}^{\prime\prime}
    &=
        \frac{1}{\sqrt{k!\left(k+\mathfrak{l}\right)!}}
        \begin{pmatrix}\mathfrak{K}\\u\end{pmatrix}
        \begin{pmatrix}\mathfrak{K}-u\\v\end{pmatrix}
        \left(k+\mathfrak{l}\right)^{\underline{v}}
        \left(1 + \mathfrak{u}\mathfrak{I}_{sv}\right)
\label{appendix:elements:precalc-coeffs-2}
,\\
\mathfrak{I}_{sv}
    &=
        \frac
        { \mathfrak{H}_{sv}^{\prime} }
        { \left(k+\mathfrak{l}+2-v\right)\left(k+\mathfrak{l}+1-v\right) }
            -
        1
\label{appendix:elements:precalc-coeffs-3}
,\\
\mathfrak{H}_{sv}^{\prime}
    &=
        \sum
        _{t=0}
        ^{\text{Min}\left[v,3+\delta_{s\bar{1}}\right]}
        v^{t}
        \mathfrak{G}_{st}
\label{appendix:elements:precalc-coeffs-4}
,\\
\mathfrak{G}_{st}
    &=
        g_{kl,s\left(1+\delta_{s\bar{1}}-t\right)}
        \frac
        {
            \sqrt
            {
                \left(k+1+\delta_{s\bar{1}}-t\right)!
                \left(k+\mathfrak{l}+2-t\right)!
            }
        }
        {
            \sqrt{k!\left(k+\mathfrak{l}\right)!}
        }
.
\label{appendix:elements:precalc-coeffs-5}
\end{align}
Please note that, in the interests of space and clarity, we did not make
explicit the fact that $\mathfrak{A}^{\prime\prime}$ and the various basis
indices and parameters within the expression should be indexed by the state
identifiers $0$ or $1$. 

It is generally less computationally intense to use the Kronecker delta to
reduce the above double summation to a single summation to implement, rather
than implementing a double summation with many zero terms. To accomplish this we
let $v_{0} \rightarrow \mathfrak{K}_{0}-v$ and $v_{1} \rightarrow
w-\mathfrak{K}_{0}+v$. Our expression becomes
\begin{equation}
\mathfrak{A}_{\left(su\right)_{01}w}^{\prime}
    =
        \sum
        _{v = \text{Max}\left[0,\mathfrak{K}_{0}+u_{1}-w\right]}
        ^{\text{Min}\left[
            \mathfrak{K}_{0}-u_{0},\mathfrak{K}_{0}+\mathfrak{K}_{1}-w
        \right]}
        \mathfrak{A}
        _{\left(su\right)_{0}\left(\mathfrak{K}_{0}-v\right)}
        ^{\prime\prime}
        \mathfrak{A}
        _{
            \left(su\right)_{1}
            \left(w-\mathfrak{K}_{0}+v\right)
        }
        ^{\prime\prime}
\label{appendix:elements:precalc-coeffs-1}
.
\end{equation}

Earlier we found some sequences of values we needed to compute to generate our matrix elements.
\begin{align}
\tilde{\bm{\mathcal{G}}}_{i\left(su\right)_{01}}
    &\coloneqq
        \mathcal{G}
        _{
            \mathfrak{h}_{\mathfrak{d}_{i}\left(s\right)_{01}};
            \left(u\right)_{01}
            \left\{\Delta\tilde{\mathcal{Z}}_{i\mathcal{G}}\right\}
        }
        ^{\left(W\right)}
\tag{\ref{definition:G-array}}
,\\
\tilde{\bm{\mathcal{H}}}_{i\left(s\right)_{01}w}
    &\coloneqq
        \mathcal{H}_{
            \left(
                \mathfrak{h}_{\mathfrak{d}_{i}\left(s\right)_{01}}
                +Ww
            \right)
            \left\{\tilde{\mathcal{Z}}_{i\mathcal{H}}\right\}
        }^{\left(W\right)}
\tag{\ref{definition:H-array}}
.
\end{align}
We note that $\mathcal{G}$ and $\mathcal{H}$ above will simplify to less
complicated forms depending on the values of $\tilde{\mathcal{Z}}$ which
themselves depend on the functions $\tilde{\alpha}$, $\beta$ and constant $W$.
If both $\delta_{W1}\tilde{\alpha} \equiv 0$ and $\delta_{W2}\beta \equiv 0$
identically, then $\mathcal{Z}_{\mathcal{G}} \equiv 1$. We also see another
simplification when $\Delta \equiv 0$ identically as well. There is a third form
when both conditions apply. To wit, we have
\begin{align}
\tilde{\mathcal{G}}_{
    L;\left(u\right)_{01}
    \left\{\Delta\left(1\right)\right\}
}^{\left(W\right)}
    &= 
        \left(-1\right)^{u_{0}}
        \Delta^{u_{0}+u_{1}}
    % && 
    %     \delta_{W1}\tilde{\alpha}_{\left\{\Omega\right\}}
    %         +
    %     \delta_{W2}\beta_{\left\{\Omega\right\}} = 0
,\\
\tilde{\mathcal{G}}_{
    L;\left(u\right)_{01}
    \left\{\left(0\right)\tilde{\mathcal{Z}}_{\mathcal{G}}\right\}
}^{\left(W\right)}
    &=
        \left(1-\tilde{\mathcal{Z}}_{\mathcal{G}}\right)^{u_{0}+u_{1}}
        \tilde{\mathcal{Z}}_{\mathcal{G}}^{\frac{L}{W}}
,\\
\tilde{\mathcal{G}}_{
        L;\left(u\right)_{01}
        \left\{\left(0\right)\left(1\right)\right\}
}^{\left(W\right)}
    &=
        \delta_{0,u_{0}+u_{1}}
.
\end{align}
We can also simplify $\tilde{\mathcal{H}}$. In the case when $W = 1$ and $\beta
= 0$ we see that $\tilde{\mathcal{Z}}_{\mathcal{H}}$ becomes ill-defined.
Returning to the original analytical integration we can see that in such a case,
$\tilde{\mathcal{H}} \equiv 1$. This can also be shown as a limit to complex
infinity, which we will denote as $\infty_{\mathbb{C}}$. Namely, we see $
\lim_{\mathcal{Z}_{\mathcal{H}}\rightarrow\infty_{\mathbb{C}}}
\tilde{\mathcal{H}} = 1 $. In essence, as $\beta$ tends to become identically
zero, it is the case that $\tilde{\mathcal{H}}$ tends to become identically
unity. A similar effect emerges when $W = 2$ and $\tilde{\alpha} \equiv 0$, but
in that case $\mathcal{Z}_{\mathcal{H}} \equiv 0$. However,
$\tilde{\mathcal{H}}$ becomes undefined for that value of
$\mathcal{Z}_{\mathcal{G}}$. As before, examining the analytic continuation
versus direct recalculation with $\tilde{\alpha} = 0$ yields the same result --
as before, under these conditions $\tilde{\mathcal{H}} \equiv 1$.

These changes to the terms of the summation over $u_{0}$, $u_{1}$, and $w$
significantly alters the structure of the summation. We present the following
contractions of the main precalculated coefficients, namely
\begin{align}
\mathfrak{A}_{\mathfrak{d}\left(s\right)_{01}uw}
    &=
        \left[ \sum_{u=0}^{\mathfrak{K}} \right]_{01}
        \delta_{u,u_{0}+u_{1}}
        \left(-1\right)^{u_{0}}
        \mathfrak{A}_{\mathfrak{d}\left(su\right)_{01}w}
,\\
\mathfrak{A}_{\mathfrak{d}\left(s\right)_{01}uw}^{\prime}
    &=
        \left[ \sum_{u=0}^{\mathfrak{K}} \right]_{01}
        \delta_{u,u_{0}+u_{1}}
        \mathfrak{A}_{\mathfrak{d}\left(su\right)_{01}w}
,\\
\mathfrak{A}_{\mathfrak{d}\left(s\right)_{01}w}
    &=
        \mathfrak{A}_{\mathfrak{d}\left(s0\right)_{01}w}
.
\end{align}
In instances when $\tilde{\mathcal{H}} \equiv 1$ we find an additional
contraction along the $w$ index, which provides the following
\begin{align}
    \mathfrak{a}_{\mathfrak{d}\left(su\right)_{01}}
        &=
            \sum_{w=u}^{\mathfrak{K}_{0}+\mathfrak{K}_{1}}
            \mathfrak{A}_{\mathfrak{d}\left(su\right)_{01}w}
    ,\\
    \mathfrak{a}_{\mathfrak{d}\left(s\right)_{01}u}
        &=
            \left[ \sum_{u=0}^{\mathfrak{K}} \right]_{01}
            \sum_{w=u}^{\mathfrak{K}_{0}+\mathfrak{K}_{1}}
            \delta_{u,u_{0}+u_{1}}
            \left(-1\right)^{u_{0}}
            \mathfrak{A}_{\mathfrak{d}\left(su\right)_{01}w}
    ,\\
    \mathfrak{a}_{\left(s\right)_{01}u}^{\prime}
        &=
            \left[ \sum_{u=0}^{\mathfrak{K}} \right]_{01}
            \sum_{w=u}^{\mathfrak{K}_{0}+\mathfrak{K}_{1}}
            \delta_{u,u_{0}+u_{1}}
            \mathfrak{A}_{\mathfrak{d}\left(su\right)_{01}w}
    ,\\
    \mathfrak{a}_{\left(s\right)_{01}}
        &=
            \sum_{w=u}^{\mathfrak{K}_{0}+\mathfrak{K}_{1}}
            \mathfrak{A}_{\mathfrak{d}\left(s0\right)_{01}w}
    .
\end{align}
As before, we note that the summations to generate these coefficients feature
alternating signs, making this computations ill-posed to be done using standard
floating-point computations. Additionally, the coefficients themselves can
become quite large. These summations have problems with both overflow and
underflow error when floating point numbers are used.

To work with this form within our numerical integrator, which will use standard
floating-point numerical computation, we opted to compute the coefficients
analytically and store the results as a floating point number. The analytic
computation can be achieved and automated with a computer algebra system, and
for this work we used the SymPy Python package. We can then store the results as
floating-point numbers in a database that can be recalled as needed. More
information is available in \ref{section:precalculation}.

Such a database can be constructed at compile time, or may be referenced from
previous runs of the software. As such, we call this set of coefficients the
\emph{precalculated coefficients}. Unfortunately, the database does prove to be
quite large, making portability between systems an issue. Additionally, we find
that pulling the needed coefficients from the database into runtime does add a
non-negligible amount of time to add computation. However, this is anticipated
to be a one-time expenditure of memory access time during a computation, one
which is anticipated to scale linearly with problem size.

In light of these simplifications and considerations we can attempt to rewrite
our integral by assessing and partitioning certain expressions. Our goal with
this formulation of the matrix elements is based on minimizing repeated
computations as well as exploiting certain features of our subbases/superbasis
structure. In particular, within a subbasis, we will have $ \eta_{0} = \eta_{1}
= \hat{\eta} $ and so we will also have $\mathfrak{f} = 0$,
$\Xi_{\left\{\Omega\right\}} \equiv 1$, and quantities such as
$\hat{\eta}/\eta_{n} = 1$. Additionally, when considering matrix elements with
$q_{0} = q_{1} = z$ within the same subbasis, we find that $\vartheta_{0} \equiv
\vartheta_{1}$ which also permits further simplifications in the expressions we
formulate below.
\begin{align}
\left(\mathcal{M}\right)_{\aleph,01}^{\left(\kappa^{\prime},j\right)}
    &\coloneqq
        \mathfrak{E}^{\mathfrak{u}_{0}+\mathfrak{u}_{1}+\mathfrak{s}}
        \mathfrak{C}
        \mathcal{I}
\label{ex:first_use_of_bracket_summation}
,\\
\mathfrak{C}
    &\coloneqq
        \frac
        { \xi_{0} }
        {
            \left[
                \left(\frac{\eta}{\hat{\eta}}\right)
                ^{\frac{\mathfrak{L}}{2}+2\mathfrak{u}}
            \right]_{01}
            \left(\frac{\hat{\eta}}{\mathfrak{r}}\right)
            ^{ 2\mathfrak{u}_{0} + 2\mathfrak{u}_{1} }
        }
,\\
\mathcal{I}
    &\coloneqq
        \int_{\mathcal{S}_{0}}\diff\Omega
        \,
        \left[
            \left(
                \frac
                {\vartheta_{\left\{\Omega\right\}}}
                {\hat{\vartheta}_{\left\{\Omega\right\}}}
            \right)^{\frac{\mathfrak{L}}{2}}
            Q_{m\zeta\left\{\phi_{q}\right\}}
        \right]_{01}
        \mathcal{S}_{\left\{\Omega\right\}}
,\\
\mathcal{S}_{\left\{\Omega\right\}}
    &\coloneqq
        \left[\sum_{s=-\mathfrak{u}}^{s=+\mathfrak{u}}\right]_{01}
        \sum_{i}
        \bm{\mathcal{P}}_{\left(s\right)_{01}\left\{\Omega\right\}}
        \bm{\mathcal{X}}_{\mathfrak{d}_{i}\left\{\Omega\right\}}
        \bm{\mathcal{V}}_{i\left\{\Omega\right\}}
        \bm{\mathcal{R}}_{i\left(s\right)_{01}\left\{\Omega\right\}}
,\\
\bm{\mathcal{P}}_{\left(s\right)_{01}\left\{\Omega\right\}}
    &\coloneqq
        \left[
            \frac
            { k!\mathfrak{H}_{ml;\lambda,s}^{\left(\kappa\right)} }
            { 
                \mathfrak{K}!
                \left(
                    \frac
                    {
                        \vartheta_{\left\{\Omega\right\}}
                        \check{r}_{\left\{\Omega\right\}}
                    }
                    {\eta}
                \right)^{W\delta_{s\bar{1}}}
            }
            P_{\left(l+2s\right)m;\lambda\left\{\chi_{q}\right\}}
        \right]_{01}
,\\
\bm{\mathcal{X}}_{\mathfrak{d}\left\{\Omega\right\}}
    &\coloneqq
        \left(
            \frac
            { \hat{\eta} }
            { \mathfrak{r}\hat{\vartheta} }
        \right)^{\mathfrak{d}}
        \left( \Xi_{\left\{\Omega\right\}} \right)
        ^{ -\hat{\mathfrak{L}}_{\mathfrak{d}} }
.
\end{align}
Above we have defined the following constants,
\begin{align}
\mathfrak{K}
    &\coloneqq
        k + \mathfrak{u} + \delta_{s\bar{1}}
,\\
\hat{\mathfrak{L}}_{\mathfrak{d}}
    &\coloneqq
        l_{0} + l_{1}
        -\delta_{W1}\left(\mathfrak{u}_{0}+\mathfrak{u}_{1}\right)
        +\mathfrak{d}+3
.
\end{align}
We can articulate the various conditions that affect the summation structure of
$\bm{\mathcal{R}}$. There are three conditions that determine which form should
be used. They are
\begin{align}
    \delta_{W1}\tilde{\alpha}_{i} + \delta_{W2}\beta_{i} = 0
    \quad\quad &\Rightarrow \quad\quad
    \tilde{\mathcal{Z}}_{i\mathcal{G}} \equiv 1
\label{condition:1a}\tag{C1a}
,\\
    \eta_{0} = \eta_{1}
    \text{ and }
    \lambda_{0} = \lambda_{1}
    \text{ and }
    q_{0} = q_{1}
    \quad\quad &\Rightarrow \quad\quad
    \Delta \equiv 0
\label{condition:1b}\tag{C1b}
,\\
    \delta_{W1}\beta_{i} + \delta_{W2}\tilde{\alpha}_{i} = 0
    \quad\quad &\Rightarrow \quad\quad
    \tilde{\mathcal{H}} \equiv 1
\label{condition:2}\tag{C2}
.
\end{align}
If condition \eqref{condition:2} is not true, then the radially integrated part
of the integrand will go as
\begin{align}
\bm{\mathcal{R}}
_{i\left(s\right)_{01}\left\{\Omega\right\}}
    &=
        \left[\sum_{u=0}^{\mathfrak{K}}\right]_{01}
        \sum_{w=u}^{\mathfrak{K}_{0}+\mathfrak{K}_{1}}
        \mathfrak{A}_{\mathfrak{d}_{i}\left(su\right)_{01}w}
        \bm{\mathcal{G}}_{i\left(u\right)_{01}\left\{\Omega\right\}}
        \tilde{\bm{\mathcal{H}}}_{i\left(s\right)_{01}w\left\{\Omega\right\}}
,\\
\bm{\mathcal{R}}
_{i\left(s\right)_{01}\left\{\Omega\right\}}
^{\left(\text{\ref{condition:1a}}\right)}
    &=
        \sum_{u=0}^{\mathfrak{K}_{0}+\mathfrak{K}_{1}}
        \sum_{w=u}^{\mathfrak{K}_{0}+\mathfrak{K}_{1}}
        \mathfrak{A}_{\mathfrak{d}_{i}\left(s\right)_{01}uw}
        \bm{\mathcal{G}}_{u\left\{\Omega\right\}}
        \tilde{\bm{\mathcal{H}}}_{i\left(s\right)_{01}w\left\{\Omega\right\}}
,\\
\bm{\mathcal{R}}
_{i\left(s\right)_{01}\left\{\Omega\right\}}
^{\left(\text{\ref{condition:1b}}\right)}
    &=
        \sum_{u=0}^{\mathfrak{K}_{0}+\mathfrak{K}_{1}}
        \sum_{w=u}^{\mathfrak{K}_{0}+\mathfrak{K}_{1}}
        \mathfrak{A}_{\mathfrak{d}_{i}\left(s\right)_{01}uw}^{\prime}
        \tilde{\bm{\mathcal{G}}}_{i\left(s\right)_{01}u\left\{\Omega\right\}}
        \tilde{\bm{\mathcal{H}}}_{i\left(s\right)_{01}w\left\{\Omega\right\}}
,\\
\bm{\mathcal{R}}
_{i\left(s\right)_{01}\left\{\Omega\right\}}
^{\left(\text{\ref{condition:1a}}+\text{\ref{condition:1b}}\right)}
    &=
        \sum_{w=0}^{\mathfrak{K}_{0}+\mathfrak{K}_{1}}
        \mathfrak{A}_{\mathfrak{d}_{i}\left(s\right)_{01}w}
        \tilde{\bm{\mathcal{H}}}_{i\left(s\right)_{01}w\left\{\Omega\right\}}
.
\end{align}
However, if condition \eqref{condition:2} is true, then the radially integrated
part of the integrand will go as
\begin{align}
\bm{\mathcal{R}}
_{i\left(s\right)_{01}\left\{\Omega\right\}}
^{\left(\text{\ref{condition:2}}\right)}
    &=
        \left[\sum_{u=0}^{\mathfrak{K}}\right]_{01}
        \mathfrak{a}_{\mathfrak{d}_{i}\left(su\right)_{01}}
        \tilde{\bm{\mathcal{G}}}_{\left(u\right)_{01}\left\{\Omega\right\}}
,\\
\bm{\mathcal{R}}
_{i\left(s\right)_{01}\left\{\Omega\right\}}
^{\left(\text{\ref{condition:2}};\text{\ref{condition:1a}}\right)}
    &=
        \sum_{u=0}^{\mathfrak{K}_{0}+\mathfrak{K}_{1}}
        \mathfrak{a}_{\mathfrak{d}_{i}\left(s\right)_{01}u}
        \bm{\mathcal{G}}_{u\left\{\Omega\right\}}
,\\
\bm{\mathcal{R}}
_{i\left(s\right)_{01}\left\{\Omega\right\}}
^{\left(\text{\ref{condition:2}};\text{\ref{condition:1b}}\right)}
    &=
        \sum_{u=0}^{\mathfrak{K}_{0}+\mathfrak{K}_{1}}
        \mathfrak{a}_{\mathfrak{d}_{i}\left(s\right)_{01}u}^{\prime}
        \tilde{\bm{\mathcal{G}}}_{i\left(s\right)_{01}u\left\{\Omega\right\}}
,\\
\bm{\mathcal{R}}
_{i\left(s\right)_{01}}
^{
    \left(\text{\ref{condition:2}};
    \text{\ref{condition:1a}}+\text{\ref{condition:1b}}\right)
}
    &= \mathfrak{a}_{\mathfrak{d}_{i}\left(s\right)_{01}}
.
\end{align}
where we employ sequences of the form
\begin{align}
\tilde{\bm{\mathcal{G}}}_{i\left(su\right)_{01}\left\{\Omega\right\}}
    &\coloneqq
        \tilde{\mathcal{G}}
        _{
            \mathfrak{h}_{\mathfrak{d}_{i}\left(s\right)_{01}};
            \left(u\right)_{01}
            \left\{
                \Delta_{\left\{\Omega\right\}}
                \tilde{\mathcal{Z}}_{i\mathcal{G}\left\{\Omega\right\}}
            \right\}
        }
        ^{\left(W\right)}
,\\
\bm{\mathcal{G}}_{u\left\{\Omega\right\}}
    &\coloneqq
        \Delta_{\left\{\Omega\right\}}^{u}
,\\
\tilde{\bm{\mathcal{G}}}_{i\left(s\right)_{01}u\left\{\Omega\right\}}
    &\coloneqq
        \tilde{\mathcal{G}}
        _{
            \mathfrak{h}_{\mathfrak{d}_{i}\left(s\right)_{01}};u
            \left\{\tilde{\mathcal{Z}}_{i\mathcal{G}\left\{\Omega\right\}}\right\}
        }
        ^{\left(W\right)}
,\\
\tilde{\bm{\mathcal{H}}}_{i\left(s\right)_{01}w\left\{\Omega\right\}}
    &\coloneqq
        \tilde{\mathcal{H}}
        _{
            \left(\mathfrak{h}_{\mathfrak{d}_{i}\left(s\right)_{01}}+Ww\right);
            \left\{
                \tilde{\mathcal{Z}}_{i\mathcal{H}\left\{\Omega\right\}}
            \right\}
        }
        ^{\left(W\right)}
.
\end{align}
The sequence values may be calculated with the following constant
\begin{equation}
\mathfrak{h}_{\mathfrak{d}\left(s\right)_{01}}
    \coloneqq
        \hat{\mathfrak{L}}_{\mathfrak{d}}
        - W\left(\delta_{s_{0}\bar{1}} + \delta_{s_{1}\bar{1}}\right)
,
\end{equation}
and with the following accessory functions.
\begin{align}
\tilde{\mathcal{G}}
_{L;\left(u\right)_{01}\left\{\Delta\tilde{\mathcal{Z}}\right\}}
^{\left(W\right)}
    &\coloneqq
        \left(1-\tilde{\mathcal{Z}}\left(1+\Delta\right)\right)^{u_{0}}
        \left(1-\tilde{\mathcal{Z}}\left(1-\Delta\right)\right)^{u_{1}}
        \tilde{\mathcal{Z}}^{\frac{L}{W}}
,\\
\tilde{\mathcal{G}}
_{L;u\left\{\tilde{\mathcal{Z}}\right\}}
^{\left(W\right)}
    &\coloneqq
        \left(1-\tilde{\mathcal{Z}}\right)^{u}
        \tilde{\mathcal{Z}}^{\frac{L}{W}}
,\\
\tilde{\mathcal{H}}_{M\left\{\tilde{\mathcal{Z}}\right\}}^{\left(W\right)}
    &\coloneqq
        \left(
            \delta_{W1}
            \frac{\sqrt{\pi}}{\Gamma\left[\frac{M+1}{2}\right]}
            \tilde{\mathcal{Z}}^{M}
            +
            \delta_{W2}
        \right)
        h_{M\left\{\tilde{\mathcal{Z}}\right\}}
,\\
\Delta_{\left\{\Omega\right\}}
    &=
        \frac
        { \mathfrak{f}+\epsilon_{\left\{\Omega\right\}} }
        { 1 + \mathfrak{f}\epsilon_{\left\{\Omega\right\}} }
,\\
\tilde{\mathcal{Z}}_{i\mathcal{G}\left\{\Omega\right\}}
&=
    \frac
    {1}
    {
        1
            +
        \delta_{W1}
        \tilde{\alpha}_{i\left\{\Omega\right\}}
        \left(\frac{ \check{r}_{\left\{\Omega\right\}} }{ \mathfrak{r} }\right)
            +
        \delta_{W2}
        \beta_{i\left\{\Omega\right\}}
        \left(
            \frac{ \check{r}_{\left\{\Omega\right\}} }{ \mathfrak{r} }
        \right)^{2}
    }
,\\
\tilde{\mathcal{Z}}_{i\mathcal{H}\left\{\Omega\right\}}
&=
    \frac
    {
        \delta_{W1}
            +
        \tilde{\alpha}_{i\left\{\Omega\right\}}
        \left(\frac{ \check{r}_{\left\{\Omega\right\}} }{ \mathfrak{r} }\right)
    }
    {
        \sqrt{
            \delta_{W2}
                +
            \beta_{i\left\{\Omega\right\}}
            \left(
                \frac{ \check{r}_{\left\{\Omega\right\}} }{ \mathfrak{r} }
            \right)^{2}
        }
    }    
.
\end{align}